\documentclass[12pt,a4paper]{article}

% ============================================================
% PACKAGES
% ============================================================
\usepackage[utf8]{inputenc}
\usepackage[T1]{fontenc}
\usepackage[french]{babel}
\usepackage{geometry}
\usepackage{graphicx}
\graphicspath{{images/}}
\usepackage{xcolor}
\usepackage{listings}
\usepackage{booktabs}
\usepackage{array}
\usepackage{hyperref}
\usepackage{fancyhdr}
\usepackage{titlesec}
\usepackage{enumitem}
\usepackage{mdframed}
\usepackage{tikz}
\usepackage{float}

% ============================================================
% MISE EN PAGE
% ============================================================
\geometry{
  top=2.5cm,
  bottom=2.5cm,
  left=2.5cm,
  right=2.5cm
}

% ============================================================
% COULEURS
% ============================================================
\definecolor{primary}{RGB}{30, 100, 180}
\definecolor{secondary}{RGB}{240, 245, 255}
\definecolor{codebg}{RGB}{245, 245, 245}
\definecolor{codeframe}{RGB}{200, 200, 200}
\definecolor{success}{RGB}{39, 174, 96}
\definecolor{warning}{RGB}{230, 126, 34}
\definecolor{darkgray}{RGB}{60, 60, 60}

% ============================================================
% EN-TETE ET PIED DE PAGE
% ============================================================
\pagestyle{fancy}
\fancyhf{}
\fancyhead[L]{\small\textcolor{primary}{\textbf{Projet M1 GIL -- Gestion de Flotte}}}
\fancyhead[R]{\small\textcolor{darkgray}{Universite de Rouen | 2025--2026}}
\fancyfoot[C]{\thepage}
\renewcommand{\headrulewidth}{0.4pt}

% ============================================================
% STYLE DES TITRES
% ============================================================
\titleformat{\section}
  {\Large\bfseries\color{primary}}
  {\thesection.}{0.8em}{}[\titlerule]

\titleformat{\subsection}
  {\large\bfseries\color{darkgray}}
  {\thesubsection.}{0.6em}{}

\titleformat{\subsubsection}
  {\normalsize\bfseries\color{darkgray}}
  {\thesubsubsection.}{0.5em}{}

% ============================================================
% STYLE DU CODE
% ============================================================
\lstset{
  backgroundcolor=\color{codebg},
  frame=single,
  rulecolor=\color{codeframe},
  basicstyle=\ttfamily\small,
  keywordstyle=\color{primary}\bfseries,
  commentstyle=\color{gray}\itshape,
  stringstyle=\color{success},
  breaklines=true,
  showstringspaces=false,
  numbers=left,
  numberstyle=\tiny\color{gray},
  numbersep=8pt,
  tabsize=2,
  captionpos=b
}

% ============================================================
% BOITE INFO
% ============================================================
\newmdenv[
  backgroundcolor=secondary,
  linecolor=primary,
  linewidth=1.5pt,
  leftline=true,
  rightline=false,
  topline=false,
  bottomline=false,
  innerleftmargin=10pt,
  innerrightmargin=10pt,
  innertopmargin=8pt,
  innerbottommargin=8pt
]{infobox}

% ============================================================
% COMMANDE PLACEHOLDER SCREENSHOT
% ============================================================
\newcommand{\screenshot}[2]{%
  \begin{figure}[H]
    \centering
    \begin{tikzpicture}
      \draw[dashed, color=gray, line width=1pt]
        (0,0) rectangle (14cm, 5cm);
      \node[color=gray] at (7cm, 2.7cm)
        {\large\textbf{[ CAPTURE D'ECRAN ]}};
      \node[color=gray] at (7cm, 2.1cm)
        {\normalsize #1};
    \end{tikzpicture}
    \caption{#2}
  \end{figure}
}

% ============================================================
% DOCUMENT
% ============================================================
\begin{document}

% ------------------------------------------------------------
% PAGE DE TITRE
% ------------------------------------------------------------
\begin{titlepage}
  \centering
  \vspace*{1cm}

  {\Large\textbf{Universite de Rouen Normandie}\\[0.3cm]
  \large Master 1 -- Genie Informatique et Logiciel}

  \vspace{1.5cm}
  \rule{\linewidth}{0.5mm}\\[0.4cm]

  {\Huge\bfseries\color{primary} Projet Gestion de Flotte de vehicules\\[0.3cm]}
  {\LARGE\bfseries Rapport de Semaine 3}\\[0.3cm]
  {\large\color{darkgray} Microservice Véhicules}

  \rule{\linewidth}{0.5mm}\\[1.5cm]

  \begin{minipage}{0.45\textwidth}
    \begin{flushleft}
      \textbf{Etudiant :}\\
      Ahcene Amouchas \\
      Florian Pepin \\[0.5cm]
      \textbf{Encadrants :}\\
      Lydia \& Luc
    \end{flushleft}
  \end{minipage}
  \hfill
  \begin{minipage}{0.45\textwidth}
    \begin{flushright}
      \textbf{Annee universitaire :}\\
      2025 -- 2026\\[0.5cm]
      \textbf{Semaine :}\\
      3 / 7
    \end{flushright}
  \end{minipage}

  \vfill
  {\large\today}
\end{titlepage}

% ------------------------------------------------------------
% TABLE DES MATIERES
% ------------------------------------------------------------
\tableofcontents
\newpage

% ------------------------------------------------------------
% INTRODUCTION
% ------------------------------------------------------------
\section{Introduction et objectifs de la semaine}

\begin{infobox}
\textbf{Objectif de la semaine 3 :} Développer le premier microservice fonctionnel du service véhicules en respectant les bonnes pratiques de
qualité logicielle : tests, couverture de code, et observabilité intégrée dès le départ.
\end{infobox}

\vspace{0.5cm}

La semaine 3 du projet marque le passage de l'infrastructure au développement métier.
L'équipe s'est concentrée sur la conception et l'implémentation du service de véhicules qui est le socle du système de gestion de flotte.

Les grandes étapes réalisées sont les suivantes :

\begin{enumerate}[leftmargin=*, label=\textbf{\arabic*.}]
  \item Développement de l'\textbf{API REST} complète (CRUD véhicules avec validation des données)
  \item Intégration d'\textbf{Apache Kafka} en tant que producer/consumer pour les événements métier
  \item Création des \textbf{résolveurs GraphQL} pour l'agrégation future entre services
  \item \textbf{Tests unitaires (>80\% couverture) et d’intégration}
  \item Instrumentation du service avec le \textbf{SDK OpenTelemetry} (traces, métriques, logs structurés)
\end{enumerate}

\section{Tests de l'API REST avec Postman}

\begin{infobox}
\textbf{Objectif :} Valider le bon fonctionnement de l'ensemble des endpoints du service
Véhicules via des tests manuels réalisés avec \textbf{Postman}, couvrant les opérations
CRUD ainsi que la gestion des assignations.
\end{infobox}

\vspace{0.5cm}

L'ensemble des routes exposées par le service Véhicules ont été testées manuellement
à l'aide de Postman. Les tests couvrent les deux ressources principales : les
\textbf{véhicules} et leurs \textbf{assignations}. Chaque requête a retourné le code
HTTP attendu, confirmant la conformité de l'API aux standards REST.

\subsection{CRUD Véhicules}

\subsubsection*{POST /vehicles — Ajout d'un véhicule}

La création d'un véhicule s'effectue via une requête \texttt{POST} sur
\texttt{/vehicles}. Le corps de la requête contient les informations du véhicule
(\texttt{plateNumber}, \texttt{brand}, \texttt{model}, \texttt{fuelType},
\texttt{mileageKm}, \texttt{vin}, \texttt{payloadCapacityKg}, \texttt{cargoVolumeM3}).
Le service retourne un \textbf{201 Created} avec l'objet créé, incluant l'UUID généré
automatiquement et le statut par défaut \texttt{available}.

\begin{figure}[H]
  \centering
  \includegraphics[width=\linewidth]{./images/postman-8.png}
  \caption{POST /vehicles — Création d'un véhicule (201 Created)}
\end{figure}

\begin{figure}[H]
  \centering
  \includegraphics[width=\linewidth]{./images/postman-9.png}
  \caption{POST /vehicles — Détails}
\end{figure}

\subsubsection*{GET /vehicles — Liste de tous les véhicules}

Une requête \texttt{GET} sur \texttt{/vehicles} retourne la liste des véhicules
non supprimés. Le mécanisme de \textbf{soft delete} est ici mis en évidence : seuls
les véhicules actifs sont retournés, même si la base de données en contient davantage.
Le serveur répond avec un \textbf{200 OK}.

\begin{figure}[H]
  \centering
  \includegraphics[width=\linewidth]{./images/postman-7.png}
  \caption{GET /vehicles — Récupération de la liste des véhicules (200 OK)}
\end{figure}

\subsubsection*{GET /vehicles/:id — Récupération par identifiant}

La récupération d'un véhicule précis s'effectue en passant son UUID en paramètre
de chemin. Le service retourne un \textbf{200 OK} avec l'intégralité des champs
du véhicule correspondant.

\begin{figure}[H]
  \centering
  \includegraphics[width=\linewidth]{./images/postman-6.png}
  \caption{GET /vehicles/:id — Récupération d'un véhicule par UUID (200 OK)}
\end{figure}

\subsubsection*{PUT /vehicles/:id — Mise à jour complète}

La mise à jour d'un véhicule via \texttt{PUT} remplace les champs modifiables
(\texttt{mileageKm}, \texttt{brand}, \texttt{model}, etc.). Le service retourne
un \textbf{200 OK} avec les données mises à jour et un \texttt{updatedAt} rafraîchi.

\begin{figure}[H]
  \centering
  \includegraphics[width=\linewidth]{./images/postman-5.png}
  \caption{PUT /vehicles/:id — Mise à jour complète d'un véhicule (200 OK)}
\end{figure}

\newpage

\subsubsection*{PATCH /vehicles/:id/status — Mise à jour du statut}

Une route dédiée permet de modifier uniquement le statut d'un véhicule (ex.
\texttt{available}, \texttt{in\_maintenance}, \texttt{on\_delivery}). Cette approche
évite d'envoyer l'intégralité du corps pour un changement ponctuel. Le service
répond avec un \textbf{200 OK} et retourne le véhicule avec son nouveau statut.

\begin{figure}[H]
  \centering
  \includegraphics[width=\linewidth]{./images/postman-4.png}
  \caption{PATCH /vehicles/:id/status — Mise à jour du statut (200 OK)}
\end{figure}

\subsubsection*{DELETE /vehicles/:id — Suppression logique}

La suppression d'un véhicule est implémentée en \textbf{soft delete} : le véhicule
n'est pas physiquement retiré de la base de données mais marqué comme supprimé.
Le service retourne un \textbf{204 No Content}, et le véhicule disparaît des
résultats des requêtes \texttt{GET}. La base de données (visible via pgAdmin)
conserve bien les deux entrées, confirmant ce comportement.

\begin{figure}[H]
  \centering
  \includegraphics[width=\linewidth]{./images/postman-3.png}
  \caption{DELETE /vehicles/:id — Suppression logique (204 No Content)}
\end{figure}

\newpage

\subsection{Gestion des assignations}

\subsubsection*{POST /vehicles/:id/assignments — Ajout d'une assignation}

L'assignation d'un conducteur à un véhicule s'effectue via une requête \texttt{POST}
sur \texttt{/vehicles/:id/assignments}. Le corps contient le \texttt{driverId},
des \texttt{notes} optionnelles et le \texttt{createdBy}. Le service répond avec
un \textbf{201 Created} et retourne l'assignation créée avec sa date de début.

\begin{figure}[H]
  \centering
  \includegraphics[width=\linewidth]{./images/postman-2.png}
  \caption{POST /vehicles/:id/assignments — Création d'une assignation (201 Created)}
\end{figure}

\subsubsection*{GET /vehicles/:id/assignments — Liste des assignations}

La récupération des assignations d'un véhicule retourne l'historique complet des
affectations de conducteurs. Le champ \texttt{endedAt} à \texttt{null} indique
une assignation en cours. Le service répond avec un \textbf{200 OK}.

\begin{figure}[H]
  \centering
  \includegraphics[width=\linewidth]{./images/postman-1.png}
  \caption{GET /vehicles/:id/assignments — Liste des assignations (200 OK)}
\end{figure}

\subsection{Bilan des tests}

\begin{center}
\begin{tabular}{|l|l|l|c|}
  \hline
  \textbf{Requête} & \textbf{Endpoint} & \textbf{Code attendu} & \textbf{Résultat} \\
  \hline
  POST   & /vehicles                       & 201 Created    & \textcolor{green!60!black}{\textbf{OK}} \\
  GET    & /vehicles                       & 200 OK         & \textcolor{green!60!black}{\textbf{OK}} \\
  GET    & /vehicles/:id                   & 200 OK         & \textcolor{green!60!black}{\textbf{OK}} \\
  PUT    & /vehicles/:id                   & 200 OK         & \textcolor{green!60!black}{\textbf{OK}} \\
  PATCH  & /vehicles/:id/status            & 200 OK         & \textcolor{green!60!black}{\textbf{OK}} \\
  DELETE & /vehicles/:id                   & 204 No Content & \textcolor{green!60!black}{\textbf{OK}} \\
  POST   & /vehicles/:id/assignments       & 201 Created    & \textcolor{green!60!black}{\textbf{OK}} \\
  GET    & /vehicles/:id/assignments       & 200 OK         & \textcolor{green!60!black}{\textbf{OK}} \\
  \hline
\end{tabular}
\end{center}

\vspace{0.3cm}
L'ensemble des 8 endpoints testés retournent les codes HTTP attendus. Le soft delete
est correctement implémenté et le filtrage des véhicules supprimés fonctionne comme
prévu.

\end{document}